\documentclass{article}
\usepackage[utf8]{inputenc}
\usepackage[T1]{fontenc}
\usepackage{geometry}
\usepackage{graphicx}
\usepackage{float}
\usepackage{hyperref}
\usepackage[english,indonesian]{babel}

\title{Optimasi Pemrosesan Template LaTeX Menggunakan Pendekatan Berbasis AI}
\author{Academic Support Agent}
\date{\today}

\begin{document}
\maketitle

\begin{abstract}

\end{abstract}

\chapter{PENDAHULUAN}


\chapter{TINJAUAN PUSTAKA}


% Original content (commented out):
% % Vorlage für Abschlussarbeiten
%
% arbeit.tex
%
% Das ist das zentrale Dokument.

% allgemeine Einstellungen
\input{einstellungen}

\begin{document}

% neue Länge als Breitenmaß für Plots einfügen
\newlength{\plotwidth}
\setlength{\plotwidth}{0.7\columnwidth}
\newlength{\plotheight}
\setlength{\plotheight}{0.66\plotwidth}
\newlength{\semiwidth}
\setlength{\semiwidth}{0.35\columnwidth}
\newlength{\semiheight}
\setlength{\semiheight}{0.66\semiwidth}

% die Titelseite
\pagenumbering{alph}
\pagestyle{empty}
\include{kapitel/titelseite}

% der Vorspann
\pagenumbering{roman}
\pagestyle{plain}
\include{kapitel/kurzfassung}
\include{kapitel/aufgabe}
\include{kapitel/erklaerung}

% die Verzeichnisse
\pagenumbering{arabic}
\pagestyle{headings}
\tableofcontents
\include{kapitel/notation}
\include{kapitel/abkuerzungen}
\listoffigures
\listoftables

% der eigentliche Inhalt
\include{kapitel/einleitung}
\include{kapitel/hauptteil}
\include{kapitel/zusammenfassung}

% die Literaturliste
\bibliographystyle{unsrtdin}	% nur Nummern, deutscher Stil nach DIN
\bibliography{literatur}

% der Anhang
\appendix
\include{kapitel/anhang}

% die Thesen
\pagestyle{empty}	
\include{kapitel/thesen}

\end{document}

% Ende der Datei


\end{document}